\newpage

\begin{center}
    {\Large\itshape Alban en sa « légende dorée »}
\end{center}

\vspace{1.2em}

\lettrine[lines=3,lhang=0.1,nindent=0em]{I}{l} y a de cela fort longtemps, vivait en son castel Alain de Nougaret, seigneur très redouté, baron d'Alban. Aujourd'hui, le château a complètement disparu. Il est mentionné en des actes notariés de 1248. On peut retrouver quelques traces du hameau de Nougaret au « Roc des Anglais » où la roche se voit taillée à angles droits. Un document du XV\textsuperscript{e} siècle situe le château de Saint Michel de Nougaret entre le ravin de Nougaret et le chemin d'Alban à Saint André. La ferme dite du « Castellàs » pourrait bien être un témoin de l'édifice aujourd'hui introuvable\ldots La réputation d'Alain de Nougaret était rien moins que bonne. On lui prêtait des pillages d'églises, des assassinats de pèlerins venus faire leurs dévotions au sanctuaire de Notre Dame d'Alban et d'ailleurs\ldots Notre Dame de Bon Secours de Villefranche, Notre Dame de la Garde par delà Saint Jean de Janes. Ce seigneur quelque peu bagarreur ne se contentait pas de détrousser les pieuses personnes que la dévotion attirait en ces lieux. Les paysans qui travaillaient ses propres terres n'étaient pas à l'abri de ses mauvais traitements. Comment de telles dépravations pouvaient-elles se concilier avec des pratiques de dévotion à l'égard de la Madone ? Bien malin qui le dira\ldots
